% ---------------------------------------------------------------------------
% Chapter 24: The Registry Office II: Top-Tagging
% Generated from the outline; edit freely, this file is now the source.
% ---------------------------------------------------------------------------
\chapter{The Registry Office II: Top-Tagging}
\chaptermark{The Registry Office II}
\label{ch:registry_office_top}

\begin{journeybox}
The registry office now has to recognise three families in one household. A mass
window and a three-prong score used to be enough; today a mass-decorrelated
network does it, and it must be calibrated on real, semileptonic $t\bar t$
events.
\end{journeybox}

\physicsline{resolved vs boosted top, CMSTopTagger, HOTVR, soft-drop mass window and $\tau_{32}$, ParticleNet-MD, mass decorrelation, scale factors from semileptonic $t\bar t$, mistag rates from QCD}

\section{When three families travel as one: resolved vs boosted top reconstruction}
\label{sec:registry_office_top:when_three_families_travel_as_one_resolv}
\todo{Write this section.}

\section{Cut-based top tagging: mass window and $\tau_{32}$}
\label{sec:registry_office_top:cut_based_top_tagging_mass_window_and_ta}
\todo{Write this section.}

\section{Machine-learned top tagging and mass decorrelation}
\label{sec:registry_office_top:machine_learned_top_tagging_and_mass_dec}
\todo{Write this section.}

\section{Calibrating the tagger: scale factors and uncertainties}
\label{sec:registry_office_top:calibrating_the_tagger_scale_factors_and}
\todo{Write this section.}

\section{Our top-jet at this stage: is it tagged?}
\label{sec:registry_office_top:our_top_jet_at_this_stage_is_it_tagged}
\todo{Numbers, plots and event display of the top-jet at this stage.}

\begin{ledger}{The Registry Office II}
  \ledgerrow{before this stage}{}{}{}{--}
  \ledgerrow{after this stage}{}{}{}{}
\end{ledger}

\section{Further reading}
\label{sec:registry_office_top:further_reading}
Official and theory references for this stage: \cite{CMS:2020poo}, \cite{Lapsien:2016zor}, \cite{Qu:2019gqs}, \cite{Kasieczka:2019dbj}.

\begin{pitfalls}
\todo{Pitfalls and FAQs for newcomers at this stage.}
\end{pitfalls}

\begin{exercises}
\item \todo{Exercise 1.}
\item \todo{Exercise 2.}
\end{exercises}
